\Askhsh{12911}{4}
\begin{ylh}{2}
    \item 4.2. Ανισώσεις 2ου Βαθμού
    \item 6.1. Η Έννοια της Συνάρτησης
\end{ylh}

Ένα δημοτικό κολυμβητήριο έχει σχήμα ορθογωνίου παραλληλογράμμου $AB\Gamma\Delta$, με διαστάσεις $15m$ και $25m$. Ο δήμος, για λόγους ασφάλειας, θέλει να κατασκευάσει γύρω από το κολυμβητήριο μια πλακοστρωμένη ζώνη με σταθερό πλάτος $x$ m ($x > 0$), όπως φαίνεται στο παρακάτω σχήμα.

\begin{center}
\begin{tikzpicture}[scale=0.2]
    % Εξωτερικό ορθογώνιο (25+2x) × (15+2x)
    \fill[gray!20] (-8,-8) rectangle (33,23);
    \draw[thick] (-8,-8) rectangle (33,23);

    % Εσωτερικό κολυμβητήριο 25×15
    \fill[cyan!15] (0,0) rectangle (25,15);
    \draw[very thick, maincolor] (0,0) rectangle (25,15);

    % Ετικέτες κορυφών
    \node[above left, maincolor] at (0,15) {A};
    \node[below left, maincolor] at (0,0) {B};
    \node[below right, maincolor] at (25,0) {$\Gamma$};
    \node[above right, maincolor] at (25,15) {$\Delta$};

    % Διαστάσεις πλάτος ζώνης
    \draw[<->, thick] (12.5,15) -- (12.5,23) node[midway, right] {$x$};
    \draw[<->, thick] (25,7.5) -- (33,7.5) node[midway, above] {$x$};

    % Κορυφές εξωτερικού
    \fill[black] (-8,-8) circle (8pt);
    \fill[black] (33,-8) circle (8pt);
    \fill[black] (-8,23) circle (8pt);
    \fill[black] (33,23) circle (8pt);
\end{tikzpicture}\captionof{figure}{Άσκηση 12911}
\end{center}
\begin{alist}
\item Να αποδείξετε ότι το εμβαδόν της ζώνης δίνεται από τη σχέση: $E(x) = 4x^2 + 80x$.
\monades{9}
\item Να βρείτε το πλάτος $x$ της ζώνης, αν αυτή έχει εμβαδόν $300 m^2$.

\monades{7}
\item Ποιο μπορεί να είναι το πλάτος της ζώνης, αν αυτή έχει εμβαδόν μικρότερο από $480 m^2$; Να αιτιολογήσετε την απάντησή σας.

\monades{9}
\end{alist}

